% Modul Belajar 10, dihasilkan oleh tools/build.py dari tools/konten/p10.py
\documentclass[a4paper]{article}
\usepackage{modulITERA}
\begin{document}
\Sampul{10}{Array Dua Dimensi dan String}{Tabel baris-kolom dengan array dua dimensi, dan teks sebagai deretan karakter}{CPMK0607 (menerapkan) dan CPMK0608 (menjelaskan) pada konsep dasar algoritma dan sistem komputer}{sekitar 3 jam belajar mandiri, ditambah 150 menit di kelas}{\texttt{02 Hands-on/P10 - Array 2D dan String - Hands-on/}}
\Bagian{Modul 10}{Peta belajar}
Ikuti urutan ini dari atas ke bawah. Setiap bagian memakai apa yang sudah dibahas di bagian sebelumnya.
\par\vspace{4pt}\noindent{\fontsize{9.5pt}{13pt}\selectfont
\begin{tabular}{@{}>{\raggedright\arraybackslash}p{0.140\textwidth}>{\raggedright\arraybackslash}p{0.840\textwidth}@{}}
\kepala{Urutan} & \kepala{Bagian}\\
1 & Tentang modul ini\\\garis
2 & Masalah pembuka: Rekap nilai kuis\\\garis
3 & Langkah 1: Array dua dimensi\\\garis
4 & Langkah 2: Total per kolom\\\garis
5 & Langkah 3: string sebagai deretan karakter\\\garis
6 & Langkah 4: Mengolah karakter\\\garis
7 & Langkah 5: Fungsi anggota string yang sering dipakai\\\garis
8 & Pesan error yang sering muncul\\\garis
9 & Latihan mandiri\\\garis
10 & Rangkuman\\\garis
11 & Cek pemahaman\\\garis
12 & Exit Ticket dan tugas\\\garis
13 & Bacaan lanjut dan persiapan minggu depan\\
\garisdasar
\end{tabular}}\par\vspace{6pt}
\Bagian{Pembuka}{Tentang modul ini}
Modul ini mendampingi Pertemuan 10. Array minggu lalu berbentuk satu baris. Minggu ini kalian menyimpan data berbentuk tabel dengan array dua dimensi, lalu melihat bahwa teks pun sebenarnya deretan karakter yang bisa diproses dengan indeks.

Isinya sama dengan slide di kelas, tetapi ditulis lebih pelan dan lebih lengkap, supaya kalian bisa mengulang sendiri di rumah tanpa harus menebak maksud slide.

\Sub{Setelah menyelesaikan modul ini, kalian bisa}
\par\vspace{4pt}\noindent{\fontsize{9.5pt}{13pt}\selectfont
\begin{tabular}{@{}>{\raggedright\arraybackslash}p{0.070\textwidth}>{\raggedright\arraybackslash}p{0.680\textwidth}>{\raggedright\arraybackslash}p{0.230\textwidth}@{}}
\kepala{No} & \kepala{Kemampuan} & \kepala{Dibahas di}\\
1 & Menulis program yang memproses data tabel dengan array dua dimensi dan mengolah teks dengan \texttt{string}: mengakses karakter, menghitung, mengubah huruf, dan memakai fungsi anggota \texttt{string} (Sub-CPMK 10.1, CPMK0607) & Langkah 1 sampai 5\\\garis
2 & Menelusuri indeks baris dan kolom yang diakses pada perulangan bersarang, serta menjelaskan pemrosesan string karakter demi karakter (Sub-CPMK 10.2, CPMK0608) & kotak Tebak dulu, Latihan 3\\
\garisdasar
\end{tabular}}\par\vspace{6pt}
Karena setiap instrumen penilaian dibagi rata antara Bagian A (menulis kode) dan Bagian B (menelusuri dan menjelaskan), yang dinilai bukan hanya program kalian jalan, tetapi juga kalian bisa menjelaskan mengapa program itu jalan.

\Sub{Cara memakai modul ini}
Setiap langkah punya pola yang sama: penjelasan konsep, kadang contoh dari kehidupan sehari-hari, lalu kode yang bisa langsung kalian ketik. Sesudah kode ada kotak \textbf{Tebak dulu}. Tulis tebakan kalian di kertas sebelum membaca \textbf{Jawaban} di bawahnya. Kebiasaan menebak lalu membuktikan inilah yang membuat konsep menempel, bukan sekadar membaca.

\begin{Penting}[Ketik, jangan salin]
Ketik ulang setiap contoh di GDB Online atau editor kalian, jangan disalin-tempel. Saat mengetik, kalian akan membuat salah ketik, membaca pesan error, dan memperbaikinya. Itu bagian dari belajar.
\end{Penting}
\Sub{Yang perlu disiapkan}
Peramban dengan akses ke onlinegdb.com, atau g++ di komputer kalian sendiri. Kertas dan alat tulis untuk menebak keluaran dan membuat trace table. Folder hands-on pertemuan ini dari repo mata kuliah.

\Bagian{Pembuka}{Masalah pembuka: Rekap nilai kuis}
Asisten menyimpan nilai 3 mahasiswa untuk 4 kuis. Dosen ingin total per mahasiswa dan rata-rata per kuis.

Dua belas variabel terlalu banyak, dan satu array satu dimensi tidak memperlihatkan mana baris mahasiswa dan mana kolom kuis.

\begin{MKeluaran}[title={Tampilan yang diinginkan}]
       K1 K2 K3 K4 | total
Rani   80 75 90 85 | 330
Budi   60 70 65 72 | 267
Siti   95 88 92 90 | 365
\end{MKeluaran}
\begin{Tebak}[Coba pikirkan]
Bagaimana kalian menunjuk nilai Budi di kuis 3? Berapa angka yang kalian butuhkan untuk menunjuknya?
\end{Tebak}
\begin{Jawaban}[Cara memecahnya]
Lihat bentuknya: Data berbentuk tabel: baris mahasiswa, kolom kuis. Tunjuk dengan dua indeks: Satu nilai ditunjuk oleh pasangan baris dan kolom. Terjemahkan: \texttt{int n[3][4]} diproses dengan \texttt{for} baris di luar, kolom di dalam.
\end{Jawaban}
\Bagian{Langkah 1}{Array dua dimensi}
Array dua dimensi adalah tabel: \texttt{int n[3][4]} punya 3 baris dan 4 kolom. Elemen ditunjuk dengan dua indeks, baris lalu kolom. \texttt{n[1][2]} adalah baris kedua, kolom ketiga.

Pemrosesan per baris memakai \texttt{for} baris di luar dan \texttt{for} kolom di dalam, persis pola perulangan bersarang dari Pertemuan 6.

\begin{MKode}{matriks.cpp}
int n[3][4] = {{80, 75, 90, 85},
               {60, 70, 65, 72},
               {95, 88, 92, 90}};
cout << n[1][2] << endl;
for (int b = 0; b < 3; b++) {
    int total = 0;
    for (int k = 0; k < 4; k++) total += n[b][k];
    cout << "Baris " << b << ": " << total << endl;
}
\end{MKode}
\begin{Tebak}[]
Sebelum menjalankan program, tulis keluarannya di kertas, karakter demi karakter.
\end{Tebak}
\begin{Jawaban}[]
\begin{MKeluaran}[]
65
Baris 0: 330
Baris 1: 267
Baris 2: 365
\end{MKeluaran}
\texttt{n[b][k]}: indeks pertama baris, kedua kolom. Keduanya mulai dari 0.
\end{Jawaban}
\Bagian{Langkah 2}{Total per kolom}
Untuk menghitung per kolom (rata-rata setiap kuis), tukar urutan perulangannya: kolom di luar, baris di dalam. Yang tidak berubah adalah urutan indeks saat mengakses: tetap \texttt{n[b][k]}.

\begin{MKode}{kolom.cpp}
int n[3][4] = {{80, 75, 90, 85},
               {60, 70, 65, 72},
               {95, 88, 92, 90}};
for (int k = 0; k < 4; k++) {
    int total = 0;
    for (int b = 0; b < 3; b++) total += n[b][k];
    cout << "K" << k + 1 << ": " << total / 3.0 << endl;
}
\end{MKode}
\begin{Tebak}[]
Sebelum menjalankan program, tulis keluarannya di kertas, karakter demi karakter.
\end{Tebak}
\begin{Jawaban}[]
\begin{MKeluaran}[]
K1: 78.3333
K2: 77.6667
K3: 82.3333
K4: 82.3333
\end{MKeluaran}
Untuk per kolom, \texttt{for} kolom di luar. Penulisan indeks tetap \texttt{n[b][k]}.
\end{Jawaban}
\Bagian{Langkah 3}{string sebagai deretan karakter}
\texttt{string} menyimpan teks sebagai deretan karakter. Setiap karakter diakses dengan indeks seperti array, dan \texttt{s.length()} memberi panjangnya. Pola yang sama dengan array berlaku: \texttt{for (int i = 0; i < panjang; i++)}.

\texttt{s.length()} bertipe tak bertanda. Membandingkannya langsung dengan \texttt{int i} memicu peringatan dengan \texttt{-Wall}, jadi ubah dulu dengan \texttt{(int) s.length()}.

\begin{MKode}{string.cpp}
string s = "Teknik Informatika";
cout << s.length() << endl;
cout << s[0] << s[7] << endl;
int spasi = 0;
for (int i = 0; i < (int) s.length(); i++) {
    if (s[i] == ' ') spasi++;
}
cout << "Spasi: " << spasi << endl;
\end{MKode}
\begin{Tebak}[]
Sebelum menjalankan program, tulis keluarannya di kertas, karakter demi karakter.
\end{Tebak}
\begin{Jawaban}[]
\begin{MKeluaran}[]
18
TI
Spasi: 1
\end{MKeluaran}
\texttt{s.length()} bertipe tak bertanda; ubah ke \texttt{int} agar perbandingan dengan \texttt{i} tidak memicu peringatan.
\end{Jawaban}
\Bagian{Langkah 4}{Mengolah karakter}
Pustaka \texttt{<cctype>} menyediakan fungsi untuk memeriksa dan mengubah satu karakter: \texttt{isalpha} (huruf), \texttt{isdigit} (angka), \texttt{isspace} (spasi), \texttt{toupper} dan \texttt{tolower}. Karena setiap karakter string bisa ditugaskan, string bisa diubah di tempat.

\begin{MKode}{karakter.cpp}
string s = "Halo ITERA 2026";
int huruf = 0, angka = 0;
for (int i = 0; i < (int) s.length(); i++) {
    if (isalpha(s[i])) huruf++;
    if (isdigit(s[i])) angka++;
    s[i] = toupper(s[i]);
}
cout << huruf << " " << angka << endl;
cout << s << endl;
\end{MKode}
\begin{Tebak}[]
Sebelum menjalankan program, tulis keluarannya di kertas, karakter demi karakter.
\end{Tebak}
\begin{Jawaban}[]
\begin{MKeluaran}[]
9 4
HALO ITERA 2026
\end{MKeluaran}
Fungsi \texttt{isalpha}, \texttt{isdigit}, \texttt{toupper}, dan \texttt{tolower} ada di \texttt{<cctype>}.
\end{Jawaban}
\Bagian{Langkah 5}{Fungsi anggota string yang sering dipakai}
Selain akses per karakter, \texttt{string} punya fungsi anggota yang praktis. Yang paling sering dipakai di mata kuliah ini adalah \texttt{length}, \texttt{substr}, \texttt{find}, dan penggabungan dengan \texttt{+}.

\par\vspace{4pt}\noindent{\fontsize{9.5pt}{13pt}\selectfont
\begin{tabular}{@{}>{\raggedright\arraybackslash}p{0.250\textwidth}>{\raggedright\arraybackslash}p{0.436\textwidth}>{\raggedright\arraybackslash}p{0.294\textwidth}@{}}
\kepala{Operasi} & \kepala{Contoh dengan s = "Praktikum"} & \kepala{Hasil}\\
panjang & \texttt{s.length()} & 9\\\garis
gabung & \texttt{s + " C++"} & \texttt{Praktikum C++}\\\garis
potong & \texttt{s.substr(0, 4)} & \texttt{Prak}\\\garis
cari & \texttt{s.find("tik")} & 4\\\garis
bandingkan & \texttt{s == "praktikum"} & false (huruf besar beda)\\
\garisdasar
\end{tabular}}\par\vspace{6pt}
\begin{Penting}[]
Membandingkan string dengan \texttt{==} membedakan huruf besar dan kecil.
\end{Penting}
\Bagian{Waspada}{Pesan error yang sering muncul}
Semua pesan di tabel ini dihasilkan g++ dengan opsi \texttt{-Wall}, sama seperti yang dipakai \texttt{cek.sh}.
\par\vspace{4pt}\noindent{\fontsize{9.5pt}{13pt}\selectfont
\begin{tabular}{@{}>{\raggedright\arraybackslash}p{0.240\textwidth}>{\raggedright\arraybackslash}p{0.270\textwidth}>{\raggedright\arraybackslash}p{0.470\textwidth}@{}}
\kepala{Kesalahan} & \kepala{Contoh} & \kepala{Pesan pertama dari g++}\\
Indeks dua dimensi pakai koma & \texttt{n[1, 2]} & \texttt{warning: left operand of comma operator has no effect}\\\garis
Kolom tidak disebut & \texttt{int n[][] = ...} & \texttt{error: declaration of 'n' as multidimensional array must have bounds for all ...}\\\garis
Kutip ganda untuk char & \texttt{s[0] = "A";} & \texttt{error: invalid conversion from 'const char*' to ...}\\\garis
Perbandingan int dengan length() & \texttt{i < s.length()} & \texttt{warning: comparison of integer expressions of different signedness: 'int' and ...}\\\garis
string tanpa nilai ke int & \texttt{int x = s;} & \texttt{error: cannot convert 'std::string' \{aka 'std::\_\_cxx11::basic\_string<char>'\} ...}\\
\garisdasar
\end{tabular}}\par\vspace{6pt}
\texttt{n[1, 2]} bukan error di semua kasus; operator koma membuatnya berarti \texttt{n[2]}. Selalu tulis dua pasang kurung siku: \texttt{n[1][2]}.

\Bagian{Latihan}{Latihan mandiri}
Latihan ini sama dengan hands-on di kelas. Semua file ada di \texttt{02 Hands-on/P10 - Array 2D dan String - Hands-on/latihan/}. Kerjakan berurutan, lalu cocokkan dengan \texttt{expected/} atau jalankan \texttt{bash cek.sh}.
\par\vspace{4pt}\noindent{\fontsize{9.5pt}{13pt}\selectfont
\begin{tabular}{@{}>{\raggedright\arraybackslash}p{0.070\textwidth}>{\raggedright\arraybackslash}p{0.360\textwidth}>{\raggedright\arraybackslash}p{0.150\textwidth}>{\raggedright\arraybackslash}p{0.400\textwidth}@{}}
\kepala{No} & \kepala{File} & \kepala{Tingkat} & \kepala{Yang dilatih}\\
1 & \texttt{handson1-matriks.cpp} & Dasar & baca dan tampilkan array 2D, total baris\\\garis
2 & \texttt{handson2-kata.cpp} & Menengah & \texttt{getline}, akses karakter, \texttt{cctype}\\\garis
3 & \texttt{handson3-transpos.cpp} & Tantangan & telusuri indeks [b][k], perbaiki tiga kesalahan\\\garis
+ & \texttt{bonus-palindrom.cpp} & Pengayaan & membandingkan karakter dari dua ujung\\
\garisdasar
\end{tabular}}\par\vspace{6pt}
\Sub{Latihan 1: handson1-matriks.cpp}
Baca nilai 3 mahasiswa × 4 kuis, tampilkan tabel dengan total per mahasiswa.

\Sub{Latihan 2: handson2-kata.cpp}
Hitung panjang, vokal, dan kata dari satu kalimat, lalu ubah ke huruf kapital.

\Sub{Latihan 3: handson3-transpos.cpp}
Transpos dan diagonal matriks 3 × 3 yang salah indeks. Tulis indeks yang dibaca di setiap langkah, perbaiki, dan jelaskan.

\Sub{Bonus: bonus-palindrom.cpp}
Periksa apakah setiap kata palindrom tanpa membedakan huruf besar-kecil.

\Sub{Latihan menelusuri}
\begin{MKode}{tebak.cpp}
int m[2][3] = {{1, 2, 3}, {4, 5, 6}};
int s = 0;
for (int k = 0; k < 3; k++) {
    s += m[1][k] - m[0][k];
}
string t = "data";
t[0] = 'D';
cout << s << " " << t + t.substr(2) << endl;
\end{MKode}
\begin{Tebak}[]
Tulis keluarannya di kertas tanpa menjalankan program. Buat trace table bila perlu.
\end{Tebak}
\begin{Jawaban}[]
\begin{MKeluaran}[]
9 Datata
\end{MKeluaran}
Setiap kolom menyumbang 3 (4−1, 5−2, 6−3), jadi \texttt{s} = 9. \texttt{t} menjadi \texttt{Data}, lalu digabung dengan \texttt{ta}.
\end{Jawaban}
\Bagian{Penutup}{Rangkuman}
\par\vspace{4pt}\noindent{\fontsize{9.5pt}{13pt}\selectfont
\begin{tabular}{@{}>{\raggedright\arraybackslash}p{0.320\textwidth}>{\raggedright\arraybackslash}p{0.660\textwidth}@{}}
\kepala{Konsep} & \kepala{Yang perlu diingat}\\
Array dua dimensi & \texttt{n[b][k]}: indeks pertama baris, kedua kolom.\\\garis
Total per kolom & Untuk per kolom, \texttt{for} kolom di luar.\\\garis
string sebagai deretan karakter & \texttt{s.length()} bertipe tak bertanda; ubah ke \texttt{int} agar perbandingan dengan \texttt{i} tidak memicu peringatan.\\\garis
Mengolah karakter & Fungsi \texttt{isalpha}, \texttt{isdigit}, \texttt{toupper}, dan \texttt{tolower} ada di \texttt{<cctype>}.\\\garis
Fungsi anggota string yang sering dipakai & Membandingkan string dengan \texttt{==} membedakan huruf besar dan kecil.\\
\garisdasar
\end{tabular}}\par\vspace{6pt}
\Bagian{Penutup}{Cek pemahaman}
Jawab tanpa menjalankan program, lalu periksa dengan menjalankannya.
\begin{enumerate}
\item Berapa banyak elemen \texttt{double t[4][6]}?
\item Untuk total per kolom, perulangan mana yang di luar?
\item Apa elemen diagonal utama matriks n × n?
\item Apa hasil \texttt{string s = "ITERA"; cout << s.substr(1, 3);}?
\item Mengapa \texttt{s[0] = "A";} error?
\end{enumerate}
\begin{Jawaban}[]
\begin{enumerate}[leftmargin=16pt]
\item 24.
\item Perulangan kolom.
\item \texttt{m[i][i]} untuk i dari 0 sampai n − 1.
\item \texttt{TER}.
\item \texttt{"A"} adalah string; satu karakter ditulis \texttt{'A'}.
\end{enumerate}
\end{Jawaban}
\Bagian{Penilaian}{Exit Ticket 10}
Dikerjakan di 25 menit terakhir kelas, sendiri, tanpa AI, internet, atau diskusi. Exit Ticket 10 adalah satu dari 14 Exit Ticket; rata-ratanya menjadi komponen berbobot 25\%.
\par\vspace{4pt}\noindent{\fontsize{9.5pt}{13pt}\selectfont
\begin{tabular}{@{}>{\raggedright\arraybackslash}p{0.080\textwidth}>{\raggedright\arraybackslash}p{0.600\textwidth}>{\raggedright\arraybackslash}p{0.100\textwidth}>{\raggedright\arraybackslash}p{0.200\textwidth}@{}}
\kepala{Soal} & \kepala{Isi} & \kepala{Poin} & \kepala{CPMK}\\
A1 & Tulis program yang menghitung rata-rata setiap kolom array 2D & 25 & CPMK0607\\\garis
A2 & Tulis program yang membalik sebuah kata dan menghitung huruf kapitalnya & 25 & CPMK0607\\\garis
B1 & Tentukan keluaran perulangan bersarang yang mengakses array 2D & 20 & CPMK0608\\\garis
B2 & Jelaskan kesalahan indeks pada potongan kode transpos & 20 & CPMK0608\\\garis
B3 & Jelaskan perbedaan \texttt{'A'} dan \texttt{"A"} beserta contohnya & 10 & CPMK0608\\
\garisdasar
\end{tabular}}\par\vspace{6pt}
\Bagian{Penilaian}{Tugas 10: Sistem Nilai Kelas}
Kumpulkan sebelum Pertemuan 11 lewat kanal pengumpulan yang diumumkan dosen.

\Sub{Bagian A: program (50 poin)}
Buat program pengelola nilai kelas dengan ketentuan berikut. Gunakan \texttt{const} untuk ukuran dan perulangan bersarang untuk operasi tabel.
\begin{enumerate}
\item Baca jumlah mahasiswa (maksimal 40) dan jumlah mata kuliah (maksimal 10)
\item Baca nama setiap mahasiswa ke array \texttt{string} dengan \texttt{getline}
\item Baca nilai setiap mahasiswa untuk setiap mata kuliah ke array dua dimensi, dengan validasi 0 sampai 100
\item Tampilkan tabel nilai dengan format rapi
\item Tampilkan rata-rata setiap mahasiswa (per baris) dan setiap mata kuliah (per kolom)
\item Tampilkan nilai tertinggi dan terendah, serta mahasiswa dengan rata-rata tertinggi
\end{enumerate}
\begin{Contoh}[Bonus, tidak wajib]
Tampilkan grafik batang sederhana rata-rata setiap mahasiswa dengan karakter \texttt{*}.
\end{Contoh}
\Sub{Bagian B: penjelasan (50 poin)}
Komentar maksimal 150 kata: gambar tabel indeks [b][k] yang diakses saat menghitung rata-rata per mata kuliah untuk 2 mahasiswa dan 3 mata kuliah, dan jelaskan perbedaan urutan perulangan untuk rata-rata per baris dan per kolom.

\Sub{Rubrik Tugas 10}
\par\vspace{4pt}\noindent{\fontsize{8.5pt}{11.5pt}\selectfont
\begin{tabular}{@{}>{\raggedright\arraybackslash}p{0.190\textwidth}>{\raggedright\arraybackslash}p{0.070\textwidth}>{\raggedright\arraybackslash}p{0.200\textwidth}>{\raggedright\arraybackslash}p{0.180\textwidth}>{\raggedright\arraybackslash}p{0.170\textwidth}>{\raggedright\arraybackslash}p{0.170\textwidth}@{}}
\kepala{Kriteria} & \kepala{Poin} & \kepala{Sangat baik 85–100} & \kepala{Baik 70–84} & \kepala{Cukup 55–69} & \kepala{Kurang <55}\\
A1 Program berjalan & 20 & tanpa error dan peringatan & ada peringatan & perlu satu perbaikan & tidak terkompilasi\\\garis
A2 Kelengkapan fitur & 15 & tabel dan semua statistik & satu fitur kurang & dua kurang & tiga atau lebih kurang\\\garis
A3 Validasi dan ketepatan & 15 & validasi tepat, hasil benar & satu salah & dua salah & sebagian besar salah\\\garis
B1 Penelusuran indeks & 30 & tabel indeks tepat, alasan jelas & satu indeks keliru & tanpa tabel & tidak ada atau disalin\\\garis
B2 Cerita error dan pengujian & 20 & pesan, sebab, perbaikan, uji & pesan tidak dikutip & hanya perbaikan & tidak ada\\
\garisdasar
\end{tabular}}\par\vspace{6pt}
Nilai kriteria = poin × skor level. Contoh: A1 di level Baik dengan skor 80 memberi 20 × 0,80 = 16 poin.

\Bagian{Penutup}{Bacaan lanjut dan persiapan minggu depan}
\begin{enumerate}
\item Deitel, P. dan Deitel, H. C++ How to Program, edisi ke-10, Bab 7 dan 21.
\item learncpp.com, Bab 5 (std::string) dan 17.
\item cppreference.com, std::basic\_string.
\item Slide \texttt{01 Slide/P10 Array Dua Dimensi dan String.pdf} dan hands-on di \texttt{02 Hands-on/P10 - Array 2D dan String - Hands-on/}.
\end{enumerate}
\begin{Penting}[Minggu depan]
Pertemuan 11 membahas searching: linear search dan binary search pada array, sejalan dengan Algoritma Pertemuan 11.
\end{Penting}
\end{document}
